\section*{Hoofdstuk \ref{ch:b2:basins-compared} --- Sedimentaire bekkens vergeleken}

\begin{solution}{exo:b2:basins-compared:1}
(a) Intracratonisch: langzame subsidentie met pulsen over honderden miljoenen jaren, zonder grote
breukvorming; de oorzaak ervan staat ter discussie. (b) Rift: mechanische verdunning van de korst tijdens de rek,
daarna thermische contractie terwijl de lithosfeer afkoelt. (c) Passieve rand: de thermische subsidentie
geërfd van het riftstadium, versterkt door de last van een dikke sedimentwig. (d) en (e)
Voorland: de lithosfeer omlaaggebogen door de last van de overschuivingsgordel, de Alpen en de Himalaya.
\end{solution}

\begin{solution}{exo:b2:basins-compared:2}
(a) $Y = 3000\times(3330 - 2300)/(3330 - 1030) = 3000\times1030/2300 = \qty{1343}{m}$.
(b) $1343 + 200 = \qty{1543}{m}$. (c) Zonder zijn last zou de sokkel op \qty{1343}{m} liggen:
$1 - 1343/3000 = \qty{55}{\percent}$ van de dikte wordt door de last ingeruimd.
\end{solution}

\begin{solution}{exo:b2:basins-compared:3}
(a) $\beta = 35/25 = 1.4$. (b) $125/1.4 = \qty{89.3}{km}$. (c) $S_i = 4334\times(1 - 1/1.4) =
\qty{1238}{m}$.
\end{solution}

\begin{solution}{exo:b2:basins-compared:4}
(a) $\tau = (1.25\times10^{5})^2/(\pi^2\times8\times10^{-7}) = \qty{1.979e15}{s}$, dat wil zeggen
$1.979\times10^{15}/3.156\times10^{13} = \qty{62.7}{Myr}$. (b) $62.7\times0.8 = \qty{50.2}{Myr}$.
(c) $62.7\times(100/125)^2 = \qty{40.1}{Myr}$. (d) $1 - e^{-1} = \qty{63}{\percent}$; $1 -
e^{-3} = \qty{95}{\percent}$.
\end{solution}

\begin{solution}{exo:b2:basins-compared:5}
(a) $0.63\,e^{-1000/1960} = 0.378$; $0.63\,e^{-2000/1960} = 0.227$. (b) $T_s = 1000 -
0.63\times1960\times(0.6004 - 0.3605) = \qty{704}{m}$. (c) $m = 2700\times704 + 1030\times296 =
\qty{2.206e6}{kg.m^{-2}}$; $\bar\rho = \qty{2206}{kg.m^{-3}}$.
\end{solution}

\begin{solution}{exo:b2:basins-compared:6}
(a) $z_2' \leftarrow 704 + 1234.8\,(1 - e^{-z_2'/1960})$: 1077, 1226, 1278, 1296, 1301, 1303,
1304; de laag is \qty{1304}{m} dik. (b) $(2700\times704 + 1030\times600)/1304 =
\qty{1932}{kg.m^{-3}}$. (c) Met \qty{304}{m}, of \qty{30}{\percent}: het water dat de compactie
eruit heeft geperst.
\end{solution}

\begin{solution}{exo:b2:basins-compared:7}
(a) $1 - e^{-100/62.7} = 0.7971$. Voor $\beta = 1.5$: $3207\times0.4135\times0.7971 =
\qty{1057}{m}$; voor $\beta = 2$: $3207\times0.6366\times0.7971 = \qty{1627}{m}$. (b) Ja: 1100 tot
\qty{1400}{m} ligt ertussen. (c) $\frac{\beta}{\pi}\sin\frac{\pi}{\beta} = 1250/(3207\times0.7971)
= 0.4890$, dus $\beta = 1.6 + 0.1\times(0.4890 - 0.4705)/(0.5204 - 0.4705) = 1.64$, binnen de
50 tot \qty{100}{\percent} rek die voor het bekken is afgeleid.
\end{solution}

\begin{solution}{exo:b2:basins-compared:8}
(a) De dichtheid van de vulling vervangt $\rho_w$ in de noemer. (b) De teller is
$125\,000\times(1/3)\times74.69 = \num{3.112e6}$; $S_i = 3.112\times10^{6}/(3184.4 - 2300) =
\qty{3519}{m}$. (c) Onder water $3.112\times10^{6}/2154.4 = \qty{1445}{m}$: het met sediment gevulde
bekken is 2.44 keer dieper.
\end{solution}

\begin{solution}{exo:b2:basins-compared:9}
(a) $2171 - 100\times3330/2300 = 2171 - 145 = \qty{2026}{m}$. (b) $Y$ verandert meter voor meter
met $W_d$: van 1971 tot \qty{2371}{m}. (c) De waterdiepte, $\pm\qty{200}{m}$, tegen
\qty{145}{m} voor het zeeniveau.
\end{solution}

\begin{solution}{exo:b2:basins-compared:10}
(a) $S_i$ heeft het teken van de haakuitdrukking in zijn teller, omdat de noemer, $3184.4 - 1030$,
positief is; ze is negatief, een opheffing, wanneer de haakuitdrukking dat is. (b) $(\rho_m - \rho_c)t_c/a =
\alpha T_1\rho_m/2$ geeft $t_c = a\,\alpha T_1\rho_m/(2(\rho_m - \rho_c)) = 125\times0.04372\times
3330/1060 = \qty{17.17}{km}$. (c) $t_c = 17.17/(1 - 0.04372\times17.17/250) = \qty{17.22}{km}$.
(d) Waar de korst vóór de rek dunner is dan ongeveer \qty{17}{km}: in oceanische lithosfeer,
of in een continentale korst die al door een eerdere riftvorming is verdund.
\end{solution}

\begin{solution}{exo:b2:basins-compared:11}
(a) Vóór de rek bevat de kolom de lithosfeer met massa $M = \rho_c t_c(1 - \alpha T_1
t_c/(2a)) + \rho_m(a - t_c)(1 - \alpha T_1(a + t_c)/(2a))$; erna $\rho_w S_i + M/\beta +
\rho_m(1 - \alpha T_1)(a - S_i - a/\beta)$. Gelijkstellen geeft $S_i[\rho_m(1 - \alpha T_1) - \rho_w] =
(1 - 1/\beta)[\rho_m(1 - \alpha T_1)a - M]$, en uitwerken van $M$ geeft $\rho_m(1 - \alpha T_1)a - M =
(\rho_m - \rho_c)t_c(1 - \alpha T_1 t_c/(2a)) - \rho_m\alpha T_1 a/2$: de formule. (b) Het eindprofiel is $T_1 z/a$; het beginprofiel ligt erboven
met $T_1 z(\beta - 1)/a$ boven $a/\beta$ en met $T_1(1 - z/a)$ daaronder. De integralen zijn $T_1
a(\beta - 1)/(2\beta^2)$ en $T_1 a(\beta - 1)^2/(2\beta^2)$, met als som $T_1 a(1 - 1/\beta)/2$.
De contractie, $\alpha$ maal dit, met water belast, geeft de uitdrukking. (c) $3.28\times10^{-5}
\times3330\times1333\times125\,000\times0.5/(2\times2300) = \qty{1978}{m}$, tegen $3207\times
2/\pi = \qty{2042}{m}$: de eerste term overschat de uiteindelijke subsidentie met \qty{64}{m},
\qty{3}{\percent}.
\end{solution}

\begin{solution}{exo:b2:basins-compared:12}
(a) $3207\times0.6366\times(1 - e^{-150/62.7}) = 2042\times0.9086 = \qty{1855}{m}$. (b) $1855\times
2300/(3330 - 2400) = \qty{4588}{m}$. (c) 6.9 tot \qty{13.8}{km}. (d) Het model laat verschillende
dingen weg die het bekken dieper zouden kunnen maken: een mantellithosfeer die sterker is uitgerekt dan de korst, een regionale
(flexurele) respons op de last van het sediment, een tijdelijke opheffing van de rand tijdens de riftvorming
gevolgd door daling, of een sedimentaanvoer die zo groot is dat het bekken overvuld raakte en op het plat
tot boven zeeniveau werd belast.
\end{solution}

\begin{solution}{pb:b2:basins-compared:1}
\textbf{1.} $0.49\,e^{-4420/3704} = 0.149$ aan de top; $0.49\,e^{-6030/3704} = 0.096$ aan de basis.
\textbf{2.} $T_s = 1610 - 0.49\times3704\times(0.3032 - 0.1963) = \qty{1416}{m}$.
\textbf{3.} Schalie: $1110 - 0.63\times1960\times(e^{-3310/1960} - e^{-4420/1960}) = 1110 - 1234.8
\times(0.1847 - 0.1049) = \qty{1011}{m}$.
\textbf{4.} Kleiig zand: $2680\times1004 + 1030\times696 = \qty{3.408e6}{kg.m^{-2}}$; krijt:
$2710\times1409 + 1030\times201 = \qty{4.025e6}{kg.m^{-2}}$; schalie: $2700\times1011 + 1030\times99 =
\qty{2.832e6}{kg.m^{-2}}$; zand: $2650\times1416 + 1030\times194 = \qty{3.952e6}{kg.m^{-2}}$. Totaal
\qty{1.4217e7}{kg.m^{-2}}, $\bar\rho_s = 14\,216\,910/6030 = \qty{2358}{kg.m^{-3}}$.
\textbf{5.} $z_2' \leftarrow 1416 + 1815\,(1 - e^{-z_2'/3704})$: 1993, 2171, 2221, 2235, 2238, 2239,
2239: \qty{2239}{m}.
\textbf{6.} $(2650\times1416 + 1030\times823)/2239 = \qty{2055}{kg.m^{-3}}$.
\textbf{7.} $Y = 2239\times(3330 - 2055)/2300 + 200 = 1241 + 200 = \qty{1441}{m}$.
\textbf{8.} $Y = 6030\times(3330 - 2358)/2300 + 100 = 2548 + 100 = \qty{2648}{m}$. De eenheden worden
op hun huidige diepten gemeten: ze zijn al zo gecompacteerd als ze vandaag zijn.
\textbf{9.} $(2171 - 1441)/50 = \qty{14.6}{m}$ per miljoen jaar van 150 tot \qty{100}{Ma};
$(2453 - 2171)/40 = \qty{7.05}{m}$ van 100 tot \qty{60}{Ma}; $(2648 - 2453)/60 = \qty{3.25}{m}$
sindsdien: de daling vertraagt.
\textbf{10.} $6030 + 100 = \qty{6130}{m}$, 2.31 keer de tektonische subsidentie.
\textbf{11.} $(6130 - 2648)/6130 = \qty{57}{\percent}$.
\textbf{12.} $2171 - 100\times3330/2300 = \qty{2026}{m}$.
\textbf{13.} Het punt aan het eind van de riftvorming, \qty{150}{Ma}: $S_i = \qty{1441}{m}$, als er tijdens de
\qty{15}{Myr} van de riftvorming geen thermische subsidentie plaatsvond.
\textbf{14.} $\beta = 1/(1 - 1441/4334) = 1/0.6675 = 1.50$.
\textbf{15.} $35/1.50 = \qty{23.3}{km}$ korst en $125/1.50 = \qty{83.3}{km}$ lithosfeer.
\textbf{16.} $(1.50/\pi)\sin(\pi/1.50) = 0.4775\times0.8660 = 0.4135$; $3207\times0.4135 =
\qty{1326}{m}$.
\textbf{17.} $S_i = 4334/3 = \qty{1445}{m}$. Op \qty{100}{Ma}, $t = \qty{50}{Myr}$: $1326\times(1 -
e^{-50/62.7}) = \qty{729}{m}$, $Y = \qty{2173}{m}$; op \qty{60}{Ma}: $1326\times0.7620 =
\qty{1010}{m}$, $Y = \qty{2455}{m}$; vandaag: $1326\times0.9086 = \qty{1205}{m}$, $Y = \qty{2650}{m}$.
\textbf{18.} Het model ligt op elk punt \qty{2}{m} dieper: 2173 tegen 2171, 2455 tegen 2453,
2650 tegen 2648.
\textbf{19.} Vroegere waterdiepten die slechts tot op honderd meter of meer bekend zijn; vroegere zeeniveaus
die verschillen van het huidige; een riftvorming die \qty{15}{Myr} duurde, waarin een deel van de warmte
al verloren ging; flexuur in plaats van lokale isostasie; een mantel die meer of minder is uitgerekt dan de korst;
een korst die vóór de riftvorming niet \qty{35}{km} dik was.
\textbf{20.} $1326\times(e^{-150/62.7} - e^{-200/62.7}) = \qty{67}{m}$.
\textbf{21.} $66.6\times2300/(3330 - 2400) = \qty{165}{m}$.
\textbf{22.} $1 - e^{-150/62.7} = \qty{91}{\percent}$.
\textbf{23.} $\tau\ln2 = \qty{43.5}{Myr}$ na het einde van de riftvorming: op \qty{106.5}{Ma}.
\textbf{24.} $\beta = \textbf{1.50}$: de korst werd uitgerekt tot anderhalve keer haar breedte.
\end{solution}
